\section{Conclusion}
\label{sec:conclusion}

The implemented MoE and stacking system achieves strong ensemble results on IMDB and Amazon under the reported thesis protocol (Section~\ref{sec:thesis_results}): STACK1 reaches macro-F1 \textbf{0.8911} and \textbf{0.8746} on those corpora respectively, while simpler ensembles and several HRM-heavy stacks lag, highlighting the value of diverse experts paired with a disciplined meta-learner and the risk of weak experts diluting stacks. Challenging surface forms (Sentiment140, TweetEval Feminism) and comparatively low mean macro-F1 for transformers and HRMs in this checkout point to future work: stronger HRM pretraining and integration, transformer hyperparameter and length tuning, imbalance-aware objectives and sampling, and deployment-oriented distillation or sparse routing.

The repository snapshot complements the narrative thesis tables with regenerated dataset alignment (\(\sim\)9.5M-row merged profile), exported model and path inventories, and HRAST stem metrics for smoke testing (Appendix~\ref{app:hrast}). Framework architecture figures (Section~\ref{sec:capstone_visuals}) document architecture and milestone language for stakeholders; all quantitative claims trace to \path{output/thesis/markdown/THESIS_IMPLEMENTATION.md}, \path{output/metrics/}, and \path{output/dataset analysis/}. The codebase structure supports swapping experts, ablations, and extended gating without rewriting the full pipeline.
